\section{Discussion and conclusions}
\label{sec:discussion-and-conclusions}

\subsection{What the measurements and tests establish}
\label{sec:result-validity}

We report the inclusive charged-current $\nu_\mu$ cross section in sixteen cells of the five-observable phase space. In the two-dimensional benchmark, the OmniFold central values reproduce MINERvA's published integral to about $1\%$. A correlated shape difference of undetermined origin remains. The five-observable central values preserve the normalization and agree with the lower-dimensional results in the tested projections. The closure tests reveal residual method bias that the reported covariances do not fully account for. At nominal truth, the negative-weight background subtraction removes most of the bias of purity subtraction in the affected muon bins and, in a separate test, at high $W$ (Secs.~\ref{sec:muon-cross-section-validation} and~\ref{sec:val-negweight} and Appendix~\ref{app:negative-weight-background}). A post hoc check finds each coarse-grid generator-test rejection also visible in a matched projection, so the tests show no evidence of discrimination beyond those projections. Table~\ref{tab:validity} gives the scope of each result.

\begingroup
\setlength{\LTcapwidth}{\textwidth}
\begin{xltabular}{\textwidth}{@{}>{\footnotesize\RaggedRight\arraybackslash\hsize=0.80\hsize}X
                                 >{\footnotesize\RaggedRight\arraybackslash\hsize=1.00\hsize}X >{\footnotesize\RaggedRight\arraybackslash\hsize=0.95\hsize}X >{\footnotesize\RaggedRight\arraybackslash\hsize=1.25\hsize}X@{}}
  \caption{Reported results, supported uses, and principal validation limits. Table~\ref{tab:higher-dimensional-validation} identifies the estimators and criteria used in the higher-dimensional validation tests.}
  \label{tab:validity}\label{tab:result-status}\\
    \toprule Result & Reported quantities & Supported use & Principal limit \\
    \midrule \endfirsthead \toprule Result & Reported quantities & Supported use & Principal limit \\
    \midrule \endhead \bottomrule \endlastfoot 2D muon cross section & \lrval{q-ncells} bin values and $\CrepTwoD$, with prior-shape dependence reported separately. & Quantitative benchmark of the integral with propagated uncertainties and descriptive comparisons of shape. & Shared inputs limit the published-result comparison, and the covariance gives no calibrated significance for generator comparisons. Closure at both held-out tunes shows a joint interval shortfall (Sec.~\ref{sec:coverage-tunes}). The reported statistical term still lacks a coverage test (Appendix~\ref{sec:coverage}). \\
    \addlinespace 3D and 4D unfoldings & Central values only, with $E_{\mathrm{avail}}$ and then $q_3$ added. & Normalization and projection checks. & Joint recovery remains unestablished, and the central values alone support no significance for generator comparisons (Secs.~\ref{sec:higher-dimensional-extension} and~\ref{sec:higher-dimensional-closure}). \\
% Search: s_proj (limitation-resolver's search, carried from sec4_higherdim.tex Sec. 4.5 at ae84584de); estimator-seed magnitude criterion; L2 (computed only for 361090f9/7e4636a3, DECISION-20260920-joseph-adopts-z-cv-under-the-6.4-exception.md:36-38; OUTCOME-20260922-L2-sproj-measured-and-the-pair-is-not-code-comparable.md:3-5) @ 0d919b1a
    \addlinespace Sixteen-cell cross section & Coarse five-observable integrals and $\CrepFiveD$. & Descriptive comparison of cell integrals with generator predictions. & Model tuning or constraints would require resolving the fake-data bias and rounding-stability failures, the dependence on the unfolding algorithm, the unexplained reconstructed run-period step (Sec.~\ref{sec:muon-cross-section-covariance}), and the interval shortfall at held-out tunes (Sec.~\ref{sec:closure-interval-diagnostics}). See also Secs.~\ref{sec:higher-dimensional-prior}--\ref{sec:higher-dimensional-calibration} and~\ref{sec:sixteen-cell-projections}. \\
    \addlinespace \ResJoint{} & Decisions on ten hybrid rate-and-shape or shape-only hypotheses. & Tests of the specified predictions within the calibrated simulation model. & The tests depend on six conditions (Appendix~\ref{sec:joint5d-conditions}). Physical exclusions of generators are outside the scope of the hybrid tests (Sec.~\ref{sec:sixteen-cell-generator-comparison}). The false-rejection rate for null data processed as in the experiment remains unmeasured. The NuWro rate-and-shape decision is sensitive to the hydrocarbon target correction (Appendix~\ref{app:joint-generator-tests}). \\
\end{xltabular}
\endgroup

\subsection{Limitations of the reported results}
\label{sec:multidimensional-limitations}
\label{sec:limitations}
\label{sec:execsummary-limitations}

The sixteen-cell estimates retain part of the prior at both the GiBUU-like truth and the synthetic $q_3$ tilt, and their cell integrals depend on the unfolding algorithm (Secs.~\ref{sec:higher-dimensional-prior} and~\ref{sec:fived16-ibu}). Coarsening the grid leaves appreciable prior dependence. Rounding-level changes also move the central values beyond the tolerance of the stability test (Sec.~\ref{sec:higher-dimensional-calibration}). The missed-event correction depends on simulated acceptance. It transfers weights learned on selected events to missed events at the same truth inputs and relies on extrapolation at very low acceptance (Sec.~\ref{sec:omnifold-implementation}).

\phantomsection
\label{sec:usage}\label{sec:limitations-use}
The covariances propagate the specified variations but are too small for the deviations of held-out-tune pseudo-experiments from the truth (Sec.~\ref{sec:closure-interval-diagnostics}). The scale of the run-period term and its ability to account for bias from pooling periods require independent checks. Linear projections of reported bins inherit the limitations of the input covariance. For inference on new observables, treating the unfolded events as independent in a likelihood can understate uncertainty because the fit correlates their weights.

\subsection{Implications for multidimensional measurements}
\label{sec:unfolding-outlook}
\label{sec:summary}

The main lesson for multidimensional measurements is that unbinned unfolding does not by itself overcome the information limits of the measured inputs. Agreement in the total rate, the muon-quadrant sums, and the marginal projections did not establish recovery of the joint cells. A shift of cross section between hadronic cells within a muon quadrant cancels in the quadrant sum, and in the five-axis $W$ injection test the $W$ marginal met its recovery criterion while the joint $(E_{\mathrm{avail}},W)$ cells in the peak interval fell below their targets. Adding an observable need not add detector information. The other inputs fix reconstructed $W$, so separating events in true $W$ relies on the simulated mapping from reconstructed variables to truth. The dependence on the prior and on the algorithm has identifiable mechanisms. In a study on noise-free samples, the prior retained at the GiBUU hybrid is consistent with a limit set by the weights of the missed events, and at the $q_3$ tilt with a limit set by the representation of the second, truth-level step (Sec.~\ref{sec:fived16-prr}). The sixteen cell integrals also differ from binned iterative Bayesian unfolding on the same inputs by more than the joint variations of the two results explain (Sec.~\ref{sec:fived16-ibu}).

Validation of the reported joint distribution requires demonstrating recovery and uncertainty coverage at the reporting resolution, including tests that vary correlations (Sec.~\ref{sec:closure-interval-diagnostics}). The origins of the run-period step and the 2D shape difference with the published result remain unknown. So does the origin of the excess of selected tracks above the reported $p_\parallel$ range (Section~\ref{sec:muon-cross-section-comparison}). The MINERvA Analysis Toolkit's signal convention does not account for the 2D shape difference with the published result (Sec.~\ref{sec:results-benchmarks}). The transfer of weights to events missed by the selection still needs validation.

% PASS2: The source prints the integral agreement as "about one percent". The equivalent literal 1\% is retained, along with the source's literal dimension labels, test counts, and q_3 subscript. (Check stage, 2026-11-21: the outlook no longer prints a study count, so the earlier "Three related studies" item is dropped.) The original summary/outlook label (sec:summary) is attached to Outlook. The original usage labels are attached to the applicability paragraph in the limitations subsection. The original top-level section wrapper was not included, so any label on that wrapper could not be verified. The restored xltabular compiled in a one-column article test. An optional REVTeX preprint test stopped inside xltabular at \@acoll.
